\documentclass[11pt,a4paper]{article}
\usepackage[margin=25mm]{geometry}
\usepackage{amsmath,amssymb,amsthm}
\setlength{\emergencystretch}{2em}
\newtheorem{lemma}{Lemma}
\newtheorem{theorem}[lemma]{Theorem}
\newcommand{\gen}[1]{\langle#1\rangle}
\newcommand{\Alt}{\operatorname{Alt}}
\title{A large pure alternating prefix at a later orbit cut}
\author{Human Compression project: L044 continuation}
\date{}
\begin{document}\maketitle
\begin{lemma}[An attached prefix near half its degree]
Let $P=\gen{s_0,\ldots,s_{d-1}}\le\Alt(C)$ be a transitive
string C-group on $c=|C|$ points, with all maximal omissions
intransitive and an anchor at its last endpoint. Suppose
\[
 d\ge c/2-1,\qquad d\ge5,\qquad
 \Alt(W)\le\gen{s_0,\ldots,s_4},\quad |W|\ge8,
\]
where $\Alt(W)$ is the actual natural group fixing $C\setminus W$.
Then $P=\Alt(C)$ naturally.
\end{lemma}
\begin{proof}
Suppose first that $P$ preserves a nontrivial equal-block system,
with block size $a\ge2$ and block number $b\ge2$.
The natural $\Alt(W)$ preserves that system. If a block meets
$W$ and its complement, an outside point is fixed by $\Alt(W)$,
so its block is invariant and contains all of $W$.
If no such mixed block occurs, the blocks meeting $W$ partition
$W$ and are invariant as a system under its primitive natural
alternating action. They are either singletons or one block.
Singleton blocks contradict $a\ge2$. In all cases some block
contains $W$, giving $a\ge8$.

The complete accepted actual-orbit forest and Whiston argument
for an anchored imprimitive interval gives $d\le a+b-2$.
This inequality uses the actual omission condition; it does not
assign IP to the block quotient. If $b\ge3$, then
\[
 \frac{ab}{2}-1-(a+b-2)
       =\frac{(a-2)(b-2)-2}{2}>0,
\]
contradicting $d\ge c/2-1$.
If $b=2$, all generators except $s_{d-1}$ fix the anchor point,
so preserve both blocks. Transitivity forces $s_{d-1}$ to
exchange them. Here $d\ge a-1\ge7$, so the string condition
makes $s_{d-1}$ commute with every generator of
$\gen{s_0,\ldots,s_4}$ and hence with $\Alt(W)$.
It must preserve the latter's unique nontrivial orbit $W$.
This is incompatible with exchanging the two blocks, one of
which contains $W$. Thus $P$ is primitive.

For completeness, a primitive permutation group containing a
natural alternating group on at least three points contains the
whole natural alternating group. Consider all conjugate supports
$W^g$, $g\in P$. They cover $C$. Two intersecting supports
can be merged: if they meet in at least two points, their pure
three-cycles generate the alternating group on the union.
If they meet only at $p$, choose distinct $u,v$ in the first
support other than $p$, and a three-cycle in the second taking
$p$ to $z$. Conjugating $(p\ u\ v)$ gives $(z\ u\ v)$.
This first merges $z$ into the first support, after which the
intersection has two points and the preceding argument applies.
The unions of connected components of the support-intersection
graph therefore form blocks for $P$, each of size at least three.
Primitivity gives one component, and the generated pure alternating
group is $\Alt(C)$. Since $P\le\Alt(C)$, equality follows.
\end{proof}

\begin{theorem}[The side containing the exceptional head]
Retain the complete L044 standing hypotheses and smaller-degree
alternating induction: $G\le A_n$, $n\ge12$, is transitive with
connected string C-marking, and every maximal omission is
intransitive. Assume rank failure $r>\lfloor(n-1)/2\rfloor$,
and the actual first three generators are an extendable source
Table~2(2) or (3) head on $D=\{1,\ldots,6\}$.
Let $j\ge5$ have a unique eligible crossing, and suppose its
omission has two orbits $C,F$ with actual supported transitive
sides
\[
       P=G_{<j}\text{ on }C,\qquad Q=G_{>j}\text{ on }F,
       \qquad D\subseteq C.
\]
The suffix is allowed to be trivial when $F$ is a singleton.
Then $P=\Alt(C)$ naturally.
\end{theorem}
\begin{proof}
The accepted critical-head reduction gives, with $E=\Omega\setminus D$,
\[
 \tau=\rho_3=(6\ b)\delta,\quad b\in E,\quad
 \operatorname{supp}(\delta)\subseteq E\setminus\{b\},
 \quad G_{\ge3}=\Alt(E\cup\{6\}).
\]
This uses its full retained primitive, anchored and small-degree
bounds and explicit induction. It does not assert that
$G_{\ge4}$ is the full alternating group on $E$.
Put $\sigma=\rho_4$. The group $G_{\ge5}$ fixes $6,b$.
If $\sigma$ fixed $b$, the transitive group $G_{\ge3}$ would
preserve $\{6,b\}$. Thus $x=b^\sigma\ne b$.
It lies outside $D$, since $\sigma$ fixes $D$ pointwise.
Conjugating $\tau$ by the transitive head and multiplying pairs
of conjugates cancels $\delta$ and supplies pure
$\Alt(D\cup\{b\})$ in $G_{\le3}$.
Conjugation by $\sigma$ then gives
\[
                  \Alt(W)\le G_{\le4},\qquad
                  W=D\cup\{b,x\},\quad |W|=8.
\]
As $j\ge5$, this group lies in $P$, so $W\subseteq C$.

The unique crossing supplies an anchor at the last endpoint of
$P$ and, when nontrivial, at the first endpoint of $Q$.
Indeed all other generators on that side commute with the
crossing involution and fix its other endpoint. Every local
maximal omission on either side is intransitive: a transitive
one, together with the entire opposite side and the crossing,
would make the corresponding ambient omission transitive.
Each side has connected string diagram, inherited from the
full connected string.

The complete paid attached-side bounds give
$\operatorname{rank}(Q)\le |F|/2$.
For degrees below 12, this follows from the complete Table~2
primitive list and the small anchored imprimitive theorem;
all listed exceptions, including the five-point dihedral and
six-point $A_5$ markings, obey this weak bound.
For larger degrees use the proper primitive and anchored
imprimitive bounds, or the explicit smaller-degree alternating
induction for a full natural alternating side. A singleton has
rank zero. No unanchored blanket bound is used.
Since $r=j+\operatorname{rank}(Q)+1\ge n/2$, it follows that
$j\ge |C|/2-1$. Apply the lemma to $P$, with $d=j$.
\end{proof}

\paragraph{Source route and exact limit.}
This handles the head-containing side at a later unshared orbit
cut in the exceptional continuation of source Propositions
5.28 and 5.32 (physical pages 36--38 and 44--45).
It replaces an unresolved group-identification alternative
by an actual natural alternating subgroup under the displayed
supported-side hypotheses. The source's cut strategy and the
accepted actual-orbit forest/Whiston bound are retained; the
large pure-subgroup block obstruction is the present argument.
The opposite dihedral or $A_5$ side, crossings with shared labels,
and the full source conclusion are not discarded. Every inherited
proof and required finite input remains paid. No new census,
formalization, original or question credit is asserted.
\end{document}
